% chapters/ch04_interrupts_timers_dma.tex — Chapter 4: Interrupts, Timers & DMA
% Scope (PLAN.md): ISR rules (what's forbidden inside, why short), latency &
% jitter, nesting & priority, shared-data problem & solutions, race conditions,
% timer/PWM frequency-duty math (worked numericals), watchdog design,
% DMA vs polling vs interrupt tradeoffs.
% Budget 40 pp, mix 22 Q&A / 14 code / 4 concept.
% 13 host snippets verified gcc -Wall -Wextra -std=c11 (zero warnings); one
% // target-only ISR skeleton syntax-checked against device-register stubs.

\chapter{Interrupts, Timers \& DMA}
\label{ch:interrupts-timers-dma}

\begin{partnote}
\textbf{Scope} --- ISR rules (what's forbidden, why short) $\rightarrow$
latency \& jitter $\rightarrow$ nesting \& priority $\rightarrow$ the
shared-data problem and its solutions $\rightarrow$ race conditions
$\rightarrow$ timer/PWM frequency--duty arithmetic (worked) $\rightarrow$
watchdog design $\rightarrow$ DMA vs polling vs interrupt. This chapter turns
the Chapter 3 NVIC/SysTick/timer hardware into correct, deterministic
real-time behaviour --- and it is where most embedded bugs actually live.
\end{partnote}

Interrupts are where ``it works on the bench'' meets ``it fails once an hour in
the field.'' Every question here is really one question: \emph{how do you share
state and time between asynchronous contexts without losing data, missing a
deadline, or corrupting a variable?} That is the daily reality of flight-control
firmware and exactly the class of HSI timing defect on the resume.

%=====================================================================
\section{Primer}
%=====================================================================

\subsection{What an ISR may and may not do}

An interrupt service routine runs in handler mode, asynchronously, often at the
worst possible moment for the code it interrupts. The rules follow from that:
keep it \textbf{short} (it stole the CPU from the main flow and may be blocking
lower-priority interrupts); do \textbf{no blocking} (no busy-wait on another
peripheral, no \code{malloc}, no long loops); do the minimum --- clear the
interrupt flag, grab the data, set a flag or push to a queue --- and defer the
heavy work to the main loop or a task. Calling non-reentrant library functions
(\code{printf}, the C heap, anything with hidden \code{static} state) from an
ISR is a classic corruption source. The mental model: an ISR \emph{captures and
hands off}; it does not \emph{process}.

\subsection{Latency, jitter, and priority}

\textbf{Interrupt latency} is the time from the hardware event to the first
useful instruction of your handler: NVIC entry (a fixed ~12 cycles on
Cortex-M, with tail-chaining/late-arrival from Chapter 3), plus any time spent
\emph{masked} by a higher-priority handler or a critical section.
\textbf{Jitter} is the variation in that latency from one occurrence to the
next --- and for a control loop, jitter is often worse than latency, because a
control law assumes a fixed sample period. The two levers are
\textbf{priority} (urgent sources get lower priority numbers so they pre-empt)
and \textbf{discipline} (short ISRs and short critical sections, so nothing
delays a pre-emption). Nesting lets a high-priority interrupt pre-empt a
low-priority handler; designing the priority map \emph{is} the real-time design.

\subsection{The shared-data problem}

The defining hazard: a variable touched by both an ISR and the main loop (or by
two interrupts). A non-atomic read-modify-write (\code{counter++}), a torn
multi-byte access, or a flag set out of order with the data it guards all
produce intermittent, hard-to-reproduce bugs. The toolkit, weakest to
strongest: \code{volatile} for \emph{visibility} (necessary, never sufficient);
a single-writer lock-free structure (the SPSC ring buffer) so neither side
needs a lock; an atomic single-word write (\code{BSRR}, bit-band, C11
\code{\_Atomic}); or a \textbf{critical section} (briefly mask interrupts) for a
true read-modify-write. The art is choosing the lightest tool that closes the
race --- a ring buffer for streaming data, a critical section for a shared
counter, a barrier-ordered flag for a handoff.

\subsection{Timers, PWM, and watchdogs}

A hardware timer is a counter clocked from the tree (Chapter 3): a
\textbf{prescaler} (\code{PSC}) divides the clock, an \textbf{auto-reload}
(\code{ARR}) sets the period, and a \textbf{compare} (\code{CCR}) sets a PWM
duty cycle. The arithmetic is exact and gets asked numerically:
\[
f_{\text{update}}=\frac{f_{\text{tim}}}{(\text{PSC}+1)(\text{ARR}+1)},\qquad
\text{duty}=\frac{\text{CCR}}{\text{ARR}+1}.
\]
A \textbf{watchdog} is a timer you must periodically ``kick'' or it resets the
MCU --- the last line of defence against a hung firmware. An independent
watchdog (IWDG) runs from its own low-speed clock so it fires even if the main
clock dies; a window watchdog (WWDG) also faults if you kick \emph{too early},
catching a runaway loop. Kicking it from a low-priority context (not a timer
ISR that would fire even if the main loop is hung) is the design subtlety.

%=====================================================================
\section{Q\&A Bank}
%=====================================================================

\tierbanner{Tier 1 --- Screening (phone / HR-technical)}%
{Two to four sentences. Interrupt fundamentals --- the gateway to any embedded
role.}

\question{Q1.1}{What is an interrupt and why use one instead of polling?}
An interrupt is a hardware signal that diverts the CPU to a handler when an
event occurs, so the processor doesn't spin waiting. It saves power and CPU
time and gives bounded response to asynchronous events --- a byte arriving, a
timer expiring --- versus polling, which wastes cycles and can miss fast
events.

\question{Q1.2}{What is the cardinal rule for an ISR?}
Keep it short and non-blocking: clear the flag, grab the data, signal the main
loop, and return. Long work, busy-waits, \code{malloc}, and non-reentrant
calls (like \code{printf}) don't belong in an ISR --- they inflate latency for
every other interrupt and risk corruption.

\question{Q1.3}{What does \code{volatile} do for an ISR-shared variable, and
what does it not do?}
It forces the compiler to actually read/write the variable each time
(visibility), so the main loop sees the ISR's updates. It does \emph{not}
provide atomicity or ordering --- a \code{counter++} is still a tearable
read-modify-write. (Chapters 1 and 3.)

\question{Q1.4}{What is interrupt latency?}
The delay from the triggering event to the first useful instruction of the
handler --- hardware entry time plus any time the interrupt is masked by a
higher-priority handler or a critical section. Lower is better; \emph{bounded}
is essential for real-time.

\question{Q1.5}{What is jitter and why does it matter for control?}
Jitter is the cycle-to-cycle variation in latency or sample timing. A control
loop (a PID) assumes a fixed period; jitter makes the effective sample time
wander, degrading stability and accuracy --- so for control, low jitter often
matters more than low average latency.

\question{Q1.6}{What is a watchdog timer?}
A timer that resets the MCU unless the firmware periodically refreshes
(``kicks'') it. It recovers a hung or runaway system automatically --- the last
line of defence when software stops making progress.

\question{Q1.7}{Timer prescaler vs auto-reload --- what does each set?}
The prescaler divides the input clock to set the counter's tick rate; the
auto-reload sets how many ticks make one period (and thus the update/overflow
frequency). Together they give
$f = f_{\text{tim}} / ((\text{PSC}+1)(\text{ARR}+1))$.

\question{Q1.8}{What is PWM and what sets the duty cycle?}
Pulse-width modulation produces a square wave of fixed frequency whose
high-time fraction (duty) encodes a value --- motor speed, LED brightness,
servo position. The compare register (\code{CCR}) relative to the reload
(\code{ARR}) sets the duty: $\text{duty}=\text{CCR}/(\text{ARR}+1)$.

\question{Q1.9}{Polling vs interrupt vs DMA in one line each?}
Polling: CPU repeatedly checks a flag (simple, wasteful). Interrupt: hardware
notifies the CPU per event (responsive, but per-event overhead). DMA: a
dedicated engine moves data without the CPU, interrupting only when a whole
block is done (best for high-throughput streams).

\question{Q1.10}{Why must you clear the interrupt flag inside the handler?}
Because most peripheral interrupts stay asserted until acknowledged; if you
don't clear the source flag, the handler re-fires immediately and the CPU is
stuck in an interrupt storm. Clearing it (and reading any required register) is
the handshake that says ``serviced.''

\tierbanner{Tier 2 --- Core technical (panel round)}%
{A short paragraph. The panel wants the mechanisms and the trade-offs.}

\question{Q2.1}{How do you safely share a counter between an ISR and the main
loop?}
For a simple counter the main loop reads while the ISR increments. The
increment \code{counter++} is a read-modify-write, so if the \emph{main loop}
also writes it, wrap that access in a short critical section (mask interrupts,
update, restore \code{PRIMASK}). If only the ISR writes and the main loop only
reads, you mainly need \code{volatile} plus an atomic read of a multi-byte
value --- e.g.\ the double-read-until-stable pattern (Coding Gym) on a core
without a 32-bit atomic load. The general rule: identify the single writer if
you can, because single-writer designs often need no lock at all.

\question{Q2.2}{Why is a lock-free SPSC ring buffer the canonical ISR-to-main
hand-off?}
Because with one producer (the ISR) writing only \code{head} and one consumer
(the main loop) writing only \code{tail}, neither side modifies the other's
index, so no lock is needed --- the ISR never blocks, which is exactly what an
ISR must not do. The producer writes the data \emph{then} advances \code{head}
(publish-after-write); the consumer reads \code{tail}'s slot then advances
\code{tail}. A power-of-two size lets the wrap be a cheap mask. It streams UART
or sensor bytes out of the ISR with bounded, deterministic cost.

\question{Q2.3}{What exactly is a race condition here, and name three fixes.}
A race is when correctness depends on the \emph{timing} of two contexts
touching shared state --- e.g.\ the main loop reads a 32-bit timestamp while the
ISR is mid-update, getting half-old/half-new bytes (a torn read), or both do
\code{x++} and one increment is lost. Fixes: (1) a \textbf{critical section}
(disable interrupts around the access); (2) an \textbf{atomic} operation
(single-word write, bit-band, \code{\_Atomic}, or LDREX/STREX); (3) a
\textbf{lock-free single-writer} structure that removes the shared
read-modify-write entirely. Choose the lightest that closes the specific race.

\question{Q2.4}{Walk the timer math: 84\,MHz timer clock, you want a 1\,kHz
interrupt --- give \code{PSC} and \code{ARR}.}
$f = f_{\text{tim}}/((\text{PSC}+1)(\text{ARR}+1))$. Pick \code{PSC}=83 to
divide 84\,MHz to 1\,MHz, then \code{ARR}=999 gives
$1\,\text{MHz}/1000 = 1\,\text{kHz}$. You split into prescaler + reload because
a single 16-bit reload can't span from a fast clock to a slow period; the
prescaler coarsely reduces the clock and the reload finely sets the period.
For a 50\% PWM you'd then set \code{CCR}=500.

\question{Q2.5}{IWDG vs WWDG --- when do you choose each?}
The \textbf{independent watchdog} runs from its own low-speed RC clock,
independent of the main clock, so it still fires if the PLL or main clock fails
--- the robust ``did the firmware die?'' guard, kicked on a coarse period. The
\textbf{window watchdog} must be kicked within a \emph{window} (not too late
\emph{and} not too early), so it also catches a runaway loop that kicks too
fast --- tighter supervision for code whose timing you want to police. Many
designs use IWDG for gross hangs; WWDG when the control loop's period itself
must be policed.

\question{Q2.6}{Where should you kick the watchdog, and why not from a timer
ISR?}
From a low-priority context that only runs when the system is genuinely making
progress --- typically the main loop after it has confirmed the critical tasks
ran this cycle. Kicking from a standalone timer ISR is a trap: that ISR keeps
firing even if the main loop is hung, so the watchdog is satisfied while the
application is dead --- defeating its purpose. A good pattern has each critical
task set a ``I ran'' flag and the kicker refresh the watchdog only when all
flags are set, then clears them.

\question{Q2.7}{DMA vs interrupt-per-byte for a 1\,Mbit/s UART stream ---
quantify why DMA wins.}
At 1\,Mbit/s a byte arrives roughly every 10\,$\mu$s. Interrupt-per-byte means
an interrupt entry/exit (tens of cycles) plus handler work every 10\,$\mu$s ---
on a slower core that's a large, jittery CPU tax and risks overrun if a
higher-priority ISR delays you. DMA moves each byte to a buffer with \emph{zero}
CPU involvement and interrupts only when the buffer (or half) is full ---
turning thousands of interrupts per second into a handful. The cost is setup
complexity and cache/coherency care; the win is CPU headroom and far lower
jitter for the rest of the system.

\question{Q2.8}{What is a half-transfer interrupt and why use a circular DMA
buffer?}
In circular (continuous) DMA the engine wraps to the buffer start
automatically, and it can interrupt at \emph{half} and \emph{full} transfer.
This gives ping-pong processing: while the CPU processes the first half, DMA
fills the second, then they swap --- continuous streaming with no gaps and no
per-byte interrupts. You compute how many new bytes arrived from the DMA's
remaining-count register (\code{NDTR}): position $=$ size $-$ \code{NDTR}, and
new bytes $=$ (position $-$ last) mod size (Coding Gym). It's the standard way
to receive an unknown-length serial frame.

\question{Q2.9}{How do you handle a 16-bit timer that wraps when measuring an
interval?}
Use unsigned modular arithmetic: \code{delta = (uint16\_t)(now - prev)}. Because
unsigned subtraction wraps modulo $2^{16}$, the difference is correct even
across one overflow (e.g.\ \code{prev}=0xFFF0, \code{now}=0x0010 gives 0x20=32).
This works as long as the true interval is less than the timer's full range; for
longer intervals you extend the timer in software by counting overflows in the
update ISR. \trap-adjacent: never compute the delta as a signed comparison ---
that breaks exactly at the wrap.

\question{Q2.10}{What is priority inversion at the interrupt level, and how does
it differ from the RTOS version?}
At the NVIC level, a long-running low-priority handler (or an over-long
critical section that masks interrupts) delays a high-priority interrupt that's
ready to run --- the urgent work waits behind the unimportant. It's prevented by
keeping ISRs and critical sections short and by correct priority assignment, not
by a protocol. The RTOS version (Chapter 9) is subtler: a medium task starves a
low task that holds a mutex the high task needs, fixed by priority inheritance.
Same symptom (urgent work blocked by less-urgent), different mechanism and fix.

\question{Q2.11}{Why prefer deferring work to a ``bottom half'' / task over
doing it in the ISR?}
Because the ISR runs at interrupt priority, blocking everything below it; long
processing there inflates latency and jitter system-wide. The ``top half'' (ISR)
does the time-critical minimum --- timestamp, read the hardware register before
it's overwritten, enqueue --- and a ``bottom half'' (main loop or a task) does
the parsing/maths at lower priority where it can be pre-empted. This keeps
interrupt latency tight and makes the heavy work schedulable, which is the basis
of responsive real-time design.

\tierbanner{Tier 3 --- Trap \& depth (principal-engineer level)}%
{A paragraph plus the trap. These are the bugs that cost field returns.}

\question{Q3.1}{An ISR sets \code{data} then sets \code{ready=1}; the main loop
sees \code{ready==1} but reads stale \code{data}. What happened and how do you
fix it without a lock?}
The compiler or the memory system reordered the two stores (or the reader's two
loads), so \code{ready} became visible before \code{data}. \code{volatile}
prevents caching but, by itself, does not guarantee ordering between the two
accesses on a core with a write buffer or weak ordering. The lock-free fix is a
\textbf{barrier} between them: producer writes \code{data}, issues a \code{DMB}
(or release-store), then writes \code{ready}; consumer reads \code{ready},
issues a \code{DMB}, then reads \code{data} (acquire). \trap the junior answer
``add volatile'' is already true and doesn't fix \emph{ordering}; the senior
answer separates visibility from ordering and names the barrier or
acquire/release semantics. This is precisely the HSI timing-defect family from
the resume.

\question{Q3.2}{Your control loop occasionally jitters by tens of
microseconds despite a hardware timer trigger. Where do you hunt?}
A timer fires deterministically, so the jitter is in the \emph{response}, not
the trigger: (1) the loop's interrupt is being delayed by a higher-or-equal
priority handler or a long critical section elsewhere (interrupt-level priority
inversion); (2) the handler itself has data-dependent execution time (a branch,
a variable-length loop, a cache effect on an ARM9); (3) flash wait-states or a
shared bus contended by DMA. The method: raise the control ISR's priority,
shrink every critical section, make the handler's path constant-time, and
measure with a GPIO toggle on a scope (the \code{DWT} cycle counter for
software). \trap the trap is blaming the timer; the timer is the one thing
that's exact. Jitter is a \emph{scheduling/contention} problem in everything
\emph{around} the handler --- exactly what an avionics 1\,kHz loop must
guarantee.

\question{Q3.3}{Is \code{counter++} in the main loop safe if only an ISR also
increments it? Prove your answer.}
No. \code{counter++} compiles to load, add, store. If the ISR fires between the
main loop's load and store, the ISR's increment is on the old value and is
overwritten by the main loop's store --- one count is lost (a lost-update
race). It's intermittent because it needs the interrupt to land in that
two-instruction window. The fix is to make the main-loop increment atomic
w.r.t.\ interrupts (critical section, or an atomic add). \trap candidates say
``it's just one variable, it's fine''; the depth is recognising that
\emph{atomicity is about the operation, not the variable's size} --- even a
single \code{int} increment is three steps. (If the main loop only \emph{reads},
the concern shrinks to a torn read of a wide value.)

\question{Q3.4}{You enabled a peripheral interrupt but the handler never runs
--- or runs once and hangs. Enumerate the likely causes.}
Never runs: the NVIC enable bit isn't set (peripheral interrupt enabled but not
the NVIC line), the priority is masked by \code{BASEPRI}/\code{PRIMASK}, the
vector table entry/name is wrong (handler not actually linked at that slot), or
the peripheral's own interrupt-enable bit is clear. Runs once then hangs: the
handler didn't clear the source flag, so the interrupt re-asserts forever (an
interrupt storm), or it re-enabled then immediately re-triggered. \trap the
two-level enable (peripheral \emph{and} NVIC) and the flag-clear handshake are
the most common bring-up mistakes; the disciplined check is peripheral-enable
$\rightarrow$ NVIC-enable $\rightarrow$ priority $\rightarrow$ vector name
$\rightarrow$ flag-clear in the handler.

\question{Q3.5}{Designing a watchdog strategy for a flight controller: what
makes it actually safe rather than theatre?}
It must detect real loss of progress, not merely ``something is still running.''
Safe design: kick from the lowest-priority context only after every critical
task has signalled completion this cycle (per-task ``I ran'' flags), so a hung
or starved task prevents the kick; use the independent watchdog (own clock) so a
clock failure still resets; size the timeout to the worst-case cycle plus
margin, not a guess; and on reset, record \emph{why} (watchdog vs power-on) so
the failure is diagnosable, and bring the system up in a safe state. \trap the
``theatre'' version kicks the watchdog from a timer ISR or unconditionally at
the top of the loop --- it then resets only on a total crash, missing the
common case of one task hung while the tick keeps running. Tying the kick to
\emph{evidence of progress} is the whole point.

\question{Q3.6}{When is DMA the wrong choice, or a source of subtle bugs?}
DMA is wrong when latency-per-event matters more than throughput (you still need
an interrupt to act on each item), when transfers are tiny (setup cost exceeds
the saving), or when the data must be transformed on the fly (DMA just moves
bytes). Subtle bugs: the CPU reading a DMA buffer while the engine is still
writing it (need half/full-transfer synchronisation or barriers); cache
coherency on cores with a data cache (the CPU sees stale cached data unless the
region is non-cacheable or invalidated); buffer alignment requirements; and
forgetting that \code{volatile} doesn't make the CPU\,/\,DMA view coherent.
\trap the trap is treating DMA as ``free and automatic''; it is a second bus
master, and every shared buffer between it and the CPU is a concurrency problem
with the same ordering/atomicity discipline as ISR-shared data.

%=====================================================================
\section{Coding Gym}
%=====================================================================

\begin{partnote}
These are the real interrupt-era patterns. Problems 1--13 are host-compilable
(verified under \code{gcc -Wall -Wextra -std=c11}); problem 14 is
\code{// target-only} (a device-register ISR skeleton), syntax-validated against
register stubs. The ring buffer (problem 1) is the single most-asked embedded
coding question.
\end{partnote}

%---------------------------------------------------------------
\begin{gymproblem}{Lock-free SPSC ring buffer (ISR $\rightarrow$ main)}

\textbf{Problem.} Implement a single-producer/single-consumer ring buffer so an
ISR can push bytes and the main loop can pop them, with no lock.

\textbf{Hints.}
\begin{itemize}
\item Producer writes only \code{head}; consumer writes only \code{tail}.
\item Power-of-two size $\rightarrow$ wrap with a mask. Full when advancing
  \code{head} would meet \code{tail}.
\end{itemize}

\attempt

\textbf{Solution.}
\begin{cgym}
#include <stdint.h>
#include <stdbool.h>
#include <stdio.h>

#define RB_SIZE 8u                  /* power of two */
typedef struct {
    uint8_t buf[RB_SIZE];
    volatile uint32_t head;         /* producer (ISR) writes this  */
    volatile uint32_t tail;         /* consumer (main) writes this */
} ring_t;

static bool rb_push(ring_t *r, uint8_t v) {           /* producer */
    uint32_t next = (r->head + 1u) & (RB_SIZE - 1u);
    if (next == r->tail) return false;                /* full -> drop */
    r->buf[r->head] = v;
    r->head = next;                                   /* publish AFTER the write */
    return true;
}
static bool rb_pop(ring_t *r, uint8_t *v) {           /* consumer */
    if (r->head == r->tail) return false;             /* empty */
    *v = r->buf[r->tail];
    r->tail = (r->tail + 1u) & (RB_SIZE - 1u);
    return true;
}

int main(void) {
    ring_t r = {{0}, 0, 0};
    rb_push(&r, 10); rb_push(&r, 20); rb_push(&r, 30);
    uint8_t v; int sum = 0;
    while (rb_pop(&r, &v)) sum += v;
    printf("sum=%d empty=%d\n", sum, rb_pop(&r, &v) == false);  /* sum=60 empty=1 */
    return 0;
}
\end{cgym}

The lock-free guarantee comes from \emph{single-writer-per-index}: the ISR only
advances \code{head}, the main loop only advances \code{tail}, so they never
modify the same variable and need no mutex --- critical because an ISR must
never block. \code{head} is advanced \emph{after} the data store so the consumer
never sees a slot before it's written (on a weakly-ordered core add a release
barrier there). One slot is left empty to distinguish full from empty.

\followups
\begin{enumerate}
\item \emph{``Why leave one slot empty?''} $\rightarrow$ So \code{head==tail}
  unambiguously means empty; full is ``next would equal tail.'' (Alternative:
  a separate count, but that's a shared RMW.)
\item \emph{``Make it byte-stream a UART RX ISR.''} $\rightarrow$ ISR calls
  \code{rb\_push} with the data register; main loop drains it.
\item \emph{``What changes on a multi-core / weak-memory system?''}
  $\rightarrow$ Need acquire/release barriers (or C11 atomics) on
  \code{head}/\code{tail}, not just \code{volatile}.
\end{enumerate}
\end{gymproblem}

%---------------------------------------------------------------
\begin{gymproblem}{Ping-pong double buffer for DMA}

\textbf{Problem.} Maintain two buffers so the DMA fills one while the CPU
processes the other, then swap.

\textbf{Hints.}
\begin{itemize}
\item Keep an \code{active} (DMA target) and a \code{processing} pointer; swap
  on the transfer-complete event.
\end{itemize}

\attempt

\textbf{Solution.}
\begin{cgym}
#include <stdint.h>
#include <stdio.h>

#define BUF_N 4
static uint8_t  bufA[BUF_N], bufB[BUF_N];
static uint8_t *active = bufA, *processing = bufB;

static void swap_buffers(void) {              /* call on transfer-complete */
    uint8_t *t = active; active = processing; processing = t;
}

int main(void) {
    for (int i = 0; i < BUF_N; i++) active[i] = (uint8_t)(i + 1);  /* DMA fills */
    swap_buffers();                                                /* now process */
    int sum = 0;
    for (int i = 0; i < BUF_N; i++) sum += processing[i];
    printf("sum=%d distinct=%d\n", sum, active != processing);     /* sum=10 distinct=1 */
    return 0;
}
\end{cgym}

While the DMA writes \code{active}, the CPU touches only \code{processing}, so
there is never a buffer both written by DMA and read by the CPU at once --- the
double-buffer removes the coherency race that a single shared buffer creates.
The swap is the only synchronised moment, done in the transfer-complete ISR. The
hardware version of this is circular DMA with half/full-transfer interrupts
(next problem's relative).

\followups
\begin{enumerate}
\item \emph{``How does circular DMA do this with one buffer?''} $\rightarrow$
  Half/full-transfer interrupts split one buffer into two halves --- same idea.
\item \emph{``Cache coherency concern?''} $\rightarrow$ On a cached core,
  invalidate the buffer before the CPU reads DMA-written data.
\item \emph{``What if processing outlasts the next fill?''} $\rightarrow$
  Overrun --- you need flow control or a third buffer/queue.
\end{enumerate}
\end{gymproblem}

%---------------------------------------------------------------
\begin{gymproblem}{PWM compare value for a duty cycle}

\textbf{Problem.} Given the auto-reload \code{ARR} and a desired duty percent,
compute the compare register \code{CCR}.

\textbf{Hints.}
\begin{itemize}
\item $\text{CCR} = (\text{ARR}+1)\times\text{duty}/100$.
\item Use a 64-bit intermediate to avoid overflow.
\end{itemize}

\attempt

\textbf{Solution.}
\begin{cgym}
#include <stdint.h>
#include <stdio.h>

static uint32_t pwm_ccr(uint32_t arr, uint32_t duty_pct) {
    return (uint32_t)(((uint64_t)(arr + 1u) * duty_pct) / 100u);
}

int main(void) {
    printf("ccr=%u\n", pwm_ccr(999u, 25u));   /* 25% of 1000 -> 250 */
    return 0;
}
\end{cgym}

The period spans \code{ARR+1} counts, so 25\% duty is \code{CCR}=250: the output
is high for counts 0--249 and low for 250--999. Duty resolution is therefore
\code{1/(ARR+1)} --- a larger \code{ARR} gives finer duty steps but (at fixed
timer clock) a lower PWM frequency, the fundamental resolution-vs-frequency
trade for motor and servo drive.

\followups
\begin{enumerate}
\item \emph{``Servo: 1.5\,ms pulse in a 20\,ms frame --- find \code{CCR}.''}
  $\rightarrow$ duty = 1.5/20 = 7.5\%; \code{CCR} = 0.075$\times$(ARR+1).
\item \emph{``Why not floating point on an MCU without an FPU?''}
  $\rightarrow$ Software FP is slow; integer scaling with a 64-bit temp keeps
  it exact and fast.
\item \emph{``Update duty without glitching.''} $\rightarrow$ Use the timer's
  preload (shadow) register so \code{CCR} changes at the update event.
\end{enumerate}
\end{gymproblem}

%---------------------------------------------------------------
\begin{gymproblem}{PWM frequency from the clock tree}

\textbf{Problem.} Compute the PWM/update frequency from the timer clock,
prescaler, and reload.

\textbf{Hints.}
\begin{itemize}
\item $f = f_{\text{tim}}/((\text{PSC}+1)(\text{ARR}+1))$.
\end{itemize}

\attempt

\textbf{Solution.}
\begin{cgym}
#include <stdint.h>
#include <stdio.h>

static uint32_t pwm_freq(uint32_t tim_clk, uint32_t psc, uint32_t arr) {
    return tim_clk / ((psc + 1u) * (arr + 1u));
}

int main(void) {
    printf("freq=%u\n", pwm_freq(84000000u, 0u, 999u));   /* 84 kHz */
    return 0;
}
\end{cgym}

With \code{PSC}=0 the counter runs at the full 84\,MHz and \code{ARR}=999 gives
84\,MHz/1000 = 84\,kHz. Note the resolution: 1000 steps means ~10-bit duty
control at 84\,kHz. Want 20-bit-ish resolution at audio rates? Increase
\code{ARR} (lower frequency) --- you cannot have both maximum frequency and
maximum resolution from one timer clock.

\followups
\begin{enumerate}
\item \emph{``Get exactly 50\,Hz for a servo from 84\,MHz.''} $\rightarrow$
  \code{PSC}=83 $\rightarrow$ 1\,MHz, \code{ARR}=19999 $\rightarrow$ 50\,Hz.
\item \emph{``Why might the real frequency differ slightly?''} $\rightarrow$
  Integer division truncation; pick \code{PSC}/\code{ARR} that divide evenly.
\item \emph{``Where does the timer clock's $\times2$ APB rule bite?''}
  $\rightarrow$ If the APB prescaler $\neq1$, \code{tim\_clk} is twice the APB
  clock.
\end{enumerate}
\end{gymproblem}

%---------------------------------------------------------------
\begin{gymproblem}{Switch debounce by sample integration}

\textbf{Problem.} In a periodic tick, sample a bouncing input and only declare a
stable state after several consecutive equal samples.

\textbf{Hints.}
\begin{itemize}
\item Shift each sample into a history byte; \code{0xFF} = stable high,
  \code{0x00} = stable low.
\end{itemize}

\attempt

\textbf{Solution.}
\begin{cgym}
#include <stdint.h>
#include <stdio.h>

/* call once per tick; 'stable' updates only after 8 consecutive equal samples */
static uint8_t debounce(uint8_t history, int raw, int *stable) {
    history = (uint8_t)((history << 1) | (raw ? 1u : 0u));
    if (history == 0xFFu)      *stable = 1;
    else if (history == 0x00u) *stable = 0;
    return history;
}

int main(void) {
    uint8_t h = 0; int st = 0;
    int samples[] = {1,0,1,1,1,1,1,1,1,1};   /* bounces, then settles high */
    for (size_t i = 0; i < sizeof samples/sizeof samples[0]; i++)
        h = debounce(h, samples[i], &st);
    printf("stable=%d\n", st);                /* 1 */
    return 0;
}
\end{cgym}

Sampling on a fixed timer tick (every few ms) and requiring 8 agreeing samples
rejects contact bounce without blocking delays --- far better than a
\code{delay\_ms()} in an ISR or a polling spin. The number of bits sets the
debounce time: 8 samples at a 1\,ms tick is an 8\,ms filter. This is the
interrupt-friendly, non-blocking way to read every mechanical switch.

\followups
\begin{enumerate}
\item \emph{``Detect the press edge, not the level.''} $\rightarrow$ Fire when
  \code{history} transitions \code{0x7F}$\rightarrow$\code{0xFF} (just became
  stable high).
\item \emph{``Tune the debounce time.''} $\rightarrow$ Change the tick period or
  the number of bits required.
\item \emph{``Why not an RC + Schmitt instead?''} $\rightarrow$ Hardware
  debounce works but costs parts; software integration is free and
  adjustable.
\end{enumerate}
\end{gymproblem}

%---------------------------------------------------------------
\begin{gymproblem}{Tear-free read of an ISR-updated 32-bit counter}

\textbf{Problem.} On a core without an atomic 32-bit load, read a counter that
an ISR increments, guaranteeing you never observe a half-updated value.

\textbf{Hints.}
\begin{itemize}
\item Read twice; if the two reads disagree, an update happened mid-read --- try
  again.
\end{itemize}

\attempt

\textbf{Solution.}
\begin{cgym}
#include <stdint.h>
#include <stdio.h>

static volatile uint32_t isr_ticks = 0;

static uint32_t read_ticks(void) {
    uint32_t a, b;
    do {
        a = isr_ticks;
        b = isr_ticks;
    } while (a != b);     /* re-read until two agree -> no tear */
    return a;
}

int main(void) {
    isr_ticks = 0x12345678u;
    printf("ticks=0x%08X\n", read_ticks());
    return 0;
}
\end{cgym}

On an 8/16-bit core a 32-bit load is several instructions, so an interrupt
between them can yield a torn value (high half new, low half old). Reading until
two consecutive reads agree guarantees no update occurred across the pair. The
alternative is a brief critical section around the read; the double-read avoids
masking interrupts at all, which is preferable when the read is frequent. (On a
32-bit core a single aligned 32-bit load is already atomic --- know your target.)

\followups
\begin{enumerate}
\item \emph{``When is this unnecessary?''} $\rightarrow$ When the value is
  $\le$ the core's atomic word size and naturally aligned.
\item \emph{``A 64-bit timestamp on a 32-bit core?''} $\rightarrow$ Same
  double-read, or read high/low/high and retry if high changed.
\item \emph{``Compare to a critical section.''} $\rightarrow$ Critical section
  adds latency to every interrupt; double-read only costs the reader.
\end{enumerate}
\end{gymproblem}

%---------------------------------------------------------------
\begin{gymproblem}{Software countdown timers driven by the tick}

\textbf{Problem.} Provide several software timers that the system tick
decrements, so the app can poll for expiry without more hardware timers.

\textbf{Hints.}
\begin{itemize}
\item An array of counters; the tick ISR decrements the non-zero ones; expiry
  is ``reached zero.''
\end{itemize}

\attempt

\textbf{Solution.}
\begin{cgym}
#include <stdint.h>
#include <stdbool.h>
#include <stdio.h>

#define NTIMERS 4
static uint32_t timers[NTIMERS];

static void tick_timers(void) {                 /* from the SysTick ISR */
    for (int i = 0; i < NTIMERS; i++)
        if (timers[i]) timers[i]--;
}
static void timer_set(int i, uint32_t ticks) { timers[i] = ticks; }
static bool timer_expired(int i) { return timers[i] == 0u; }

int main(void) {
    timer_set(0, 3); timer_set(1, 5);
    for (int t = 0; t < 3; t++) tick_timers();
    printf("t0=%d t1=%d\n", timer_expired(0), timer_expired(1));   /* 1 0 */
    return 0;
}
\end{cgym}

One hardware tick fans out into many cheap software timeouts --- the basis of
``do X in 50\,ms'' scheduling without dedicating a peripheral timer to each. The
ISR work is trivial (a short decrement loop), keeping it within the
short-ISR rule, and the app polls expiry at its leisure. Watch the shared-access
caveat: \code{timer\_set} from the main loop while the ISR decrements is a write
race on a 32-bit value --- atomic on a 32-bit core, otherwise wrap in a
critical section.

\followups
\begin{enumerate}
\item \emph{``Auto-reload (periodic) timers.''} $\rightarrow$ On expiry, reload
  from a stored period instead of stopping.
\item \emph{``Run a callback on expiry.''} $\rightarrow$ Store a function
  pointer per timer; the main loop calls it when expired (not the ISR).
\item \emph{``Scale to hundreds of timers.''} $\rightarrow$ A sorted
  delta-queue / timer wheel instead of scanning every tick.
\end{enumerate}
\end{gymproblem}

%---------------------------------------------------------------
\begin{gymproblem}{Interval measurement across a 16-bit timer wrap}

\textbf{Problem.} Compute elapsed ticks between two readings of a 16-bit
free-running timer that may have wrapped once.

\textbf{Hints.}
\begin{itemize}
\item Unsigned subtraction is modular --- just cast the difference back to
  16-bit.
\end{itemize}

\attempt

\textbf{Solution.}
\begin{cgym}
#include <stdint.h>
#include <stdio.h>

static uint16_t elapsed(uint16_t now, uint16_t prev) {
    return (uint16_t)(now - prev);    /* modulo-2^16 handles the wrap */
}

int main(void) {
    printf("delta=%u\n", elapsed(0x0010u, 0xFFF0u));   /* wrapped: 0x20 = 32 */
    return 0;
}
\end{cgym}

Because unsigned arithmetic wraps modulo $2^{16}$, \code{now - prev} is the
correct elapsed count even when \code{now < prev} due to a single overflow ---
no branch, no special case. The only requirement is that the true interval be
less than the timer's full span ($2^{16}$ ticks); longer intervals need a
software-extended timer that counts overflows in the update ISR. Never use a
signed comparison here --- it fails precisely at the wrap.

\followups
\begin{enumerate}
\item \emph{``Extend to intervals longer than 65535 ticks.''} $\rightarrow$
  Count overflows in the update ISR to form a 32-bit time base.
\item \emph{``Why is \code{(uint16\_t)} cast needed?''} $\rightarrow$ Integer
  promotion makes \code{now-prev} an \code{int}; the cast re-wraps it to 16
  bits.
\item \emph{``Measure a pulse width with input capture.''} $\rightarrow$ The
  timer latches the count on the edge; subtract two captures the same way.
\end{enumerate}
\end{gymproblem}

%---------------------------------------------------------------
\begin{gymproblem}{Independent-watchdog timeout calculation}

\textbf{Problem.} Compute the IWDG timeout in milliseconds from its low-speed
clock, prescaler, and reload value.

\textbf{Hints.}
\begin{itemize}
\item $t = \text{prescaler}\times(\text{reload}+1)/f_{\text{LSI}}$.
\end{itemize}

\attempt

\textbf{Solution.}
\begin{cgym}
#include <stdint.h>
#include <stdio.h>

static uint32_t iwdg_timeout_ms(uint32_t lsi_hz, uint32_t presc, uint32_t reload) {
    return (uint32_t)((uint64_t)presc * (reload + 1u) * 1000u / lsi_hz);
}

int main(void) {
    printf("timeout=%u ms\n", iwdg_timeout_ms(32000u, 32u, 1249u));   /* 1250 ms */
    return 0;
}
\end{cgym}

A 32\,kHz LSI, prescaler 32, reload 1249 gives
$32\times1250/32000 = 1.25\,\text{s}$. You size the timeout to the worst-case
loop period plus margin: too short and a legitimate slow cycle resets you; too
long and a hang goes undetected for ages. The LSI's own tolerance (it's an RC
oscillator) means you leave generous margin rather than cutting it fine.

\followups
\begin{enumerate}
\item \emph{``Pick values for a 4\,s timeout.''} $\rightarrow$ Solve
  \code{presc*(reload+1) = 4*lsi}; e.g.\ presc=128, reload=999.
\item \emph{``Why account for LSI tolerance?''} $\rightarrow$ The RC clock
  varies several percent with temperature/voltage; margin prevents spurious
  resets.
\item \emph{``Where do you kick it?''} $\rightarrow$ Low-priority context, only
  after confirming progress (Q2.6/Q3.5).
\end{enumerate}
\end{gymproblem}

%---------------------------------------------------------------
\begin{gymproblem}{Interrupt response-time budget}

\textbf{Problem.} Given the NVIC entry cycles and the handler's worst-case
cycles, compute the response time in nanoseconds at a given core clock.

\textbf{Hints.}
\begin{itemize}
\item $t = (\text{entry}+\text{work})/f_{\text{cpu}}$; use a 64-bit
  intermediate and report ns.
\end{itemize}

\attempt

\textbf{Solution.}
\begin{cgym}
#include <stdint.h>
#include <stdio.h>

static uint32_t response_ns(uint32_t entry_cyc, uint32_t work_cyc, uint32_t cpu_hz) {
    return (uint32_t)((uint64_t)(entry_cyc + work_cyc) * 1000000000u / cpu_hz);
}

int main(void) {
    printf("resp=%u ns\n", response_ns(12u, 100u, 168000000u));   /* ~666 ns */
    return 0;
}
\end{cgym}

12 cycles of NVIC entry plus a 100-cycle handler at 168\,MHz is ~666\,ns ---
\emph{if} nothing masks the interrupt. The real worst case adds the longest time
the interrupt can be blocked: a higher-priority handler's duration, or the
longest critical section anywhere in the system. That ``blocking term'' is
usually the dominant and most overlooked part of the budget, which is why short
critical sections matter so much.

\followups
\begin{enumerate}
\item \emph{``Add a 2\,$\mu$s critical section elsewhere --- new worst case?''}
  $\rightarrow$ Add 2000\,ns; it likely dwarfs the entry+work term.
\item \emph{``How do you measure entry cycles for real?''} $\rightarrow$
  Toggle a GPIO at the event and at handler entry; read on a scope, or use
  \code{DWT->CYCCNT}.
\item \emph{``Why is average latency misleading?''} $\rightarrow$ Real-time
  cares about the worst case, not the mean --- one missed deadline is a
  failure.
\end{enumerate}
\end{gymproblem}

%---------------------------------------------------------------
\begin{gymproblem}{Peak-to-peak jitter from a timestamp series}

\textbf{Problem.} Given measured inter-arrival intervals, compute the
peak-to-peak jitter.

\textbf{Hints.}
\begin{itemize}
\item Jitter = max interval $-$ min interval.
\end{itemize}

\attempt

\textbf{Solution.}
\begin{cgym}
#include <stdint.h>
#include <stddef.h>
#include <stdio.h>

static uint32_t jitter(const uint32_t *intervals, size_t n) {
    uint32_t mn = intervals[0], mx = intervals[0];
    for (size_t i = 1; i < n; i++) {
        if (intervals[i] < mn) mn = intervals[i];
        if (intervals[i] > mx) mx = intervals[i];
    }
    return mx - mn;     /* peak-to-peak */
}

int main(void) {
    uint32_t iv[] = {1000, 1002, 998, 1001, 999};
    printf("jitter=%u\n", jitter(iv, sizeof iv/sizeof iv[0]));   /* 4 */
    return 0;
}
\end{cgym}

A nominal 1000-tick period that ranges 998--1002 has 4 ticks of peak-to-peak
jitter. For a control loop you quote jitter as a fraction of the period (here
0.4\%) and check it against the loop's tolerance. You'd collect these intervals
on-target from an input-capture timer or the cycle counter, then watch how
jitter grows when DMA or higher-priority ISRs contend --- the empirical version
of Q3.2.

\followups
\begin{enumerate}
\item \emph{``Report RMS jitter instead of peak-to-peak.''} $\rightarrow$
  Standard deviation of the intervals; peak-to-peak is the safety-critical
  bound.
\item \emph{``Where do you capture these timestamps cheaply?''} $\rightarrow$
  Timer input-capture on the event edge, or \code{DWT->CYCCNT} in software.
\item \emph{``A single huge outlier --- mean hides it, what catches it?''}
  $\rightarrow$ Max / peak-to-peak; always track the worst case.
\end{enumerate}
\end{gymproblem}

%---------------------------------------------------------------
\begin{gymproblem}{Rate-limited action (do every N ticks)}

\textbf{Problem.} Run an action only every \code{N} calls --- e.g.\ a slow task
driven from a fast tick.

\textbf{Hints.}
\begin{itemize}
\item A counter that fires and resets at \code{N}.
\end{itemize}

\attempt

\textbf{Solution.}
\begin{cgym}
#include <stdint.h>
#include <stdbool.h>
#include <stdio.h>

static bool every_n(uint32_t *counter, uint32_t n) {
    if (++(*counter) >= n) { *counter = 0; return true; }   /* fire & reset */
    return false;
}

int main(void) {
    uint32_t c = 0; int fires = 0;
    for (int i = 0; i < 10; i++) if (every_n(&c, 3)) fires++;
    printf("fires=%d\n", fires);   /* fires at 3,6,9 -> 3 */
    return 0;
}
\end{cgym}

This builds a slow software clock from a fast one --- run a 1\,kHz tick but
execute a 100\,Hz task by firing every 10th tick, or blink an LED every 500\,ms
from a 1\,ms tick. It keeps all timing referenced to one hardware tick (so tasks
stay phase-coherent) instead of spreading timing across many peripherals. Each
task gets its own counter and divisor.

\followups
\begin{enumerate}
\item \emph{``Stagger two tasks so they don't run on the same tick.''}
  $\rightarrow$ Give them different initial counter offsets (phase).
\item \emph{``Why keep one tick source?''} $\rightarrow$ Phase coherence and a
  single time base simplify reasoning and reduce drift.
\item \emph{``Drift over time?''} $\rightarrow$ Integer division of one accurate
  tick has no accumulating drift, unlike independent oscillators.
\end{enumerate}
\end{gymproblem}

%---------------------------------------------------------------
\begin{gymproblem}{New bytes available from a circular DMA (\code{NDTR})}

\textbf{Problem.} A circular DMA counts down \code{NDTR} from the buffer size.
Compute how many new bytes arrived since the last check, handling wrap.

\textbf{Hints.}
\begin{itemize}
\item Current position = size $-$ \code{NDTR}; new = (position $-$ last) mod
  size.
\end{itemize}

\attempt

\textbf{Solution.}
\begin{cgym}
#include <stdint.h>
#include <stdio.h>

static uint32_t dma_new_bytes(uint32_t size, uint32_t ndtr, uint32_t last_pos,
                              uint32_t *new_pos) {
    uint32_t pos   = size - ndtr;                  /* where DMA has written to */
    uint32_t avail = (pos + size - last_pos) % size;  /* new bytes, wrap-safe  */
    *new_pos = pos;
    return avail;
}

int main(void) {
    uint32_t np = 0;
    uint32_t got = dma_new_bytes(64u, 60u, 0u, &np);   /* pos=4 -> 4 new */
    printf("new=%u pos=%u\n", got, np);
    return 0;
}
\end{cgym}

In circular receive the DMA never stops, so you poll \code{NDTR} (or act on the
half/full-transfer and idle-line interrupts) to discover how far it has written.
Position $=$ size $-$ \code{NDTR} because \code{NDTR} counts \emph{down}; the
modulo handles the wrap so a read that straddles the buffer end still returns the
right count. This is the standard idiom for receiving variable-length UART
frames over DMA without an interrupt per byte (Q2.8).

\followups
\begin{enumerate}
\item \emph{``Detect end-of-frame.''} $\rightarrow$ The UART IDLE-line
  interrupt fires on a gap; read \code{NDTR} then to bound the frame.
\item \emph{``Copy the new bytes out, handling wrap.''} $\rightarrow$ Two
  \code{memcpy}s if the region crosses the buffer end.
\item \emph{``Coherency on a cached core?''} $\rightarrow$ Invalidate the
  buffer region before reading DMA-written data.
\end{enumerate}
\end{gymproblem}

%---------------------------------------------------------------
\begin{gymproblem}{Timer-update ISR skeleton (\texttt{// target-only})}

\textbf{Problem.} Write a correct timer update-interrupt handler and the
two-level enable (peripheral interrupt + NVIC), following the short-ISR and
flag-clear rules.

\textbf{Hints.}
\begin{itemize}
\item Check then clear the update flag; do the minimum; enable both the timer's
  \code{DIER} bit and the NVIC line.
\end{itemize}

\attempt

\textbf{Solution.}
\begin{cgym}
#include <stdint.h>
// target-only: uses device register definitions (TIM2, NVIC) from the CMSIS header

static volatile uint32_t tick;

void TIM2_IRQHandler(void) {
    if (TIM2->SR & TIM_SR_UIF) {     /* update event pending? */
        TIM2->SR &= ~TIM_SR_UIF;     /* clear the flag (rc_w0) BEFORE returning */
        tick++;                      /* short work only -- defer the rest        */
    }
}

void tim2_enable_irq(void) {
    TIM2->DIER |= (1u << 0);                                 /* timer: update IRQ on */
    NVIC->ISER[TIM2_IRQn >> 5] = (1u << (TIM2_IRQn & 31));   /* NVIC: enable the line */
}
\end{cgym}

Three correctness points an interviewer checks: the \textbf{flag clear} (omit it
and the handler re-fires forever --- the Q3.4 interrupt storm); the
\textbf{two-level enable} (the peripheral's \code{DIER} \emph{and} the NVIC
\code{ISER} --- forgetting the NVIC line is the ``handler never runs'' bug); and
\textbf{short work} (increment a tick, set a flag --- never parse or print
here). On STM32 the update flag is clear-by-writing-0; on other parts it may be
read-then-write or write-1-to-clear --- always check the datasheet.

\followups
\begin{enumerate}
\item \emph{``Why clear the flag early, not late?''} $\rightarrow$ If the event
  recurs during the handler, clearing late can drop or double it; clear it once,
  deterministically.
\item \emph{``How does the main loop consume \code{tick} safely?''}
  $\rightarrow$ It's a single 32-bit value: atomic on a 32-bit core; otherwise
  the double-read of problem 6.
\item \emph{``Set this timer's priority below a control ISR.''} $\rightarrow$
  Give the control interrupt a lower priority number so it pre-empts the tick.
\end{enumerate}
\end{gymproblem}

%=====================================================================
\section{Link to Her Resume}
%=====================================================================

\begin{resumelink}
\textbf{Shared-data / timing defects:} ``The HSI timing defects I resolved at
BEL are the interrupt-shared-data family --- a flag set out of order with the
data it guards, or a torn read of a multi-byte value. I fix them by separating
visibility, ordering, and atomicity: \code{volatile} plus a barrier or a short
\code{PRIMASK} critical section, and lock-free single-writer buffers where I
can.''

\medskip
\textbf{Determinism / jitter:} ``A flight-control loop lives or dies on jitter,
not average latency. I keep the control ISR high-priority, keep critical
sections short, and measure timing with a GPIO toggle on the scope --- the same
lab instruments on my resume.''

\medskip
\textbf{DMA / streaming:} ``For high-rate telemetry I'd take a UART over circular
DMA with half/full-transfer interrupts rather than an interrupt per byte ---
it's how you free the CPU and keep jitter down, computing new bytes from
\code{NDTR}.''

\medskip
\small Maps to concrete resume lines: BEL HSI timing-defect resolution,
real-time DPRAM polling with sub-1\,ms latency, and oscilloscope-based debugging
on the TASL boards. Keep claims to what the resume states; the deep project
scripting is Chapter 15, from \code{inputs/INTAKE.md} only.
\end{resumelink}

%=====================================================================
\section{Self-Test Scorecard}
%=====================================================================

\begin{scorecard}
\begin{enumerate}
\item State the cardinal ISR rule and two things forbidden inside an ISR.
\item Latency vs jitter --- which usually matters more for a control loop, and
  why?
\item Why is a lock-free SPSC ring buffer safe without a mutex?
\item Is \code{counter++} in the main loop safe if an ISR also increments it?
  Why or why not?
\item 84\,MHz timer, want 1\,kHz --- give \code{PSC} and \code{ARR}.
\item Duty formula: with \code{ARR}=999, what \code{CCR} gives 25\%?
\item Why should you not kick the watchdog from a standalone timer ISR?
\item Compute elapsed for a 16-bit timer with \code{prev}=0xFFF0,
  \code{now}=0x0010.
\item DMA vs interrupt-per-byte for a 1\,Mbit/s stream --- one-line why DMA
  wins.
\item An ISR writes \code{data} then \code{ready=1}; the reader sees stale
  \code{data}. Name the cause and the fix.
\end{enumerate}
\end{scorecard}

\begin{answerkey}
\begin{enumerate}
\item Keep it short and non-blocking; forbidden: blocking/busy-wait,
  \code{malloc}, long loops, non-reentrant calls like \code{printf} (any two).
\item Jitter --- a control law assumes a fixed sample period, so timing
  variation degrades stability more than a constant latency does.
\item Single writer per index: the ISR writes only \code{head}, the main loop
  only \code{tail}, so they never modify the same variable.
\item No --- \code{counter++} is load/add/store; an interrupt between load and
  store loses the ISR's update (lost-update race). Make it atomic / critical
  section.
\item \code{PSC}=83 (84\,MHz$\rightarrow$1\,MHz), \code{ARR}=999
  (1\,MHz/1000 = 1\,kHz).
\item \code{CCR} = 25\% of (ARR+1) = 0.25$\times$1000 = 250.
\item Because it fires regardless of whether the main loop is alive, so it
  satisfies the watchdog while the application is hung --- defeating the point.
\item \code{(uint16\_t)(0x0010 - 0xFFF0)} = 0x20 = 32 (modular subtraction
  survives the wrap).
\item DMA moves bytes with no CPU per-byte cost, interrupting only per block ---
  thousands of ISRs/s become a few, with far less jitter.
\item The two stores (or the reader's loads) were reordered; \code{volatile}
  gives visibility but not ordering. Fix: a barrier / release-store between
  \code{data} and \code{ready} (and acquire on the reader).
\end{enumerate}
\end{answerkey}

\begin{partnote}
\emph{End of Chapter 4, and of Part A (Embedded C \& Processor Core).} Next:
Part B opens with \textbf{UART / RS232 / RS422 / RS485} --- frame anatomy, baud
math and error percentage, differential signalling, and why flight computers
use RS422 (the daily interface on the resume).
\end{partnote}
